\documentclass[11pt,a4paper]{article}
\usepackage{cmap}
\usepackage[utf8]{inputenc}
\usepackage[T1]{fontenc}
\usepackage{lmodern}
\usepackage{microtype}
\usepackage[margin=1.0in]{geometry}
\usepackage{amsmath,amssymb,amsthm,mathtools}
\usepackage{booktabs,array,tabularx}
\usepackage[svgnames]{xcolor}
\usepackage[obeyspaces,spaces]{url}
\usepackage[colorlinks=true,linkcolor=DarkBlue,citecolor=DarkBlue,urlcolor=DarkBlue]{hyperref}

\urlstyle{tt}
\def\UrlFont{\ttfamily\upshape}
\renewcommand{\arraystretch}{1.3}
\setlength{\tabcolsep}{5pt}
\sloppy
\setlength{\emergencystretch}{4em}
\hfuzz=4.0pt
\vfuzz=4.0pt

\newtheorem{theorem}{Theorem}[section]
\theoremstyle{definition}
\newtheorem{definition}[theorem]{Definition}
\theoremstyle{remark}
\newtheorem{remark}{Remark}[section]

\title{\textbf{Recombination Does Not Help:\\
a Kernel-Gated Measurement of Crossover in Program Synthesis,\\
and the Verifier Artifact That Made It Affordable}}
\author{Xavier Callens\\
\small SocrateAI Scientific Research \& Formal Verification Laboratory}
\date{Edition 2, 24 October 2026. \textbf{Not peer reviewed.}}

\begin{document}
\maketitle

\begin{abstract}
Crossover is the operator that distinguishes a genetic algorithm from a
hill-climber, and in evolutionary program synthesis its contribution is usually
asserted rather than measured. Three preceding papers in this programme asked
whether recombination is \emph{detectable} in evolved programs and refuted the
question's instrument each time. None asked the prior question: does
recombination \emph{help}?

That question needs a formal gate over hundreds of candidates, and the gate cost
one proof-assistant invocation per candidate, which put a properly powered
paired experiment out of reach. This paper reports, first, that the cost was an
artifact of our own verifier wrapper rather than of the proof assistant. A Lean~4
compile exits nonzero on a batch containing a false obligation but still prints a
per-declaration axiom line for every candidate, the failing ones carrying
\texttt{sorryAx}. Our wrapper discarded those lines the instant the exit code was
nonzero, which made a batch appear atomic and motivated three successive rounds
of isolation machinery --- a bisection, a flat fall-back measured as better than
the bisection, and a break-even law $(1-p)^b > 1/b$ derived from the cost of
isolating a failure. All three were solving a problem the wrapper had invented.
Reading the kernel's own report gives every verdict in one invocation at any
failure rate, measured at $1$ invocation for $0,1,2,4$ and $8$ failures of $8$,
with every culprit named and no false candidate accepted.

With the gate affordable, we measure crossover directly. Fitness is exact
rational agreement with a specification at sample points, computed outside the
proof assistant; acceptance is the Lean kernel, so a wrong fitness can misguide
the search but cannot manufacture an acceptance. Over $200$ paired seeds on a
calibrated target, mutation-only solved $94/200$ ($47.0\%$) against crossover's
$58/200$ ($29.0\%$), Fisher exact $p = 2.98\times10^{-4}$. Conditional on
solving, the evaluation count is unchanged (median ratio $0.959$ over $35$ pairs,
sign test $p = 1.000$). Crossover therefore does not make successful searches
faster; it makes success less likely --- the signature of a macromutation
destroying assembled structure rather than a recombination transmitting it.

We give the mechanism exactly. An expression is affine in any one subtree's
value, $\text{out} = A + Bv$, and $B$ is computable in exact rationals by one
forward and one backward pass. A donor contributing $B_{\text{donor}}v$ in its
own parent contributes $B_{\text{site}}v$ after transplant, so its contribution
is rescaled or sign-inverted whenever the coefficients differ, even though its
syntax arrives intact. We then test the repair this diagnosis suggests ---
restricting recombination to \emph{homologous} sites, those whose coefficient
sign pattern matches the donor's, by analogy with homologous recombination in
biology. \textbf{The repair fails, and its own negative control refutes the
mechanism.} Homology-restricted recombination reaches $70/200$ ($35.0\%$),
neither significantly above random crossover ($p = 0.238$) nor recovering
mutation-only, which beats it significantly ($p = 0.019$); and the
anti-homologous control, restricted to \emph{inverted} sign patterns and
predicted to be worse, instead scores higher at $79/200$ ($39.5\%$) and is
statistically indistinguishable from the treatment ($p = 0.408$) --- the outcome
we had registered in advance as meaning the signature carries nothing. A
fallback-to-mutation confound is excluded quantitatively. \textbf{Edition 2 tests the successor
edition 1 recorded as untested} --- inserting the donor \emph{rescaled} so its
contribution transfers exactly, verified exact at $1614$ of $1614$ comparisons
--- and its pre-registered falsifier fires too: $73/200$ ($36.5\%$), not
significantly above random crossover ($p = 0.136$), still significantly below
mutation-only ($p = 0.042$), and indistinguishable from the weaker sign-matching
operator ($p = 0.835$). Six independent results now agree that subtree
recombination in this representation does not help, including both repairs its
exact diagnosis recommends.

Every numeric claim is reproduced by a script in the accompanying repository, and
every reported defect of our own apparatus is recorded with the measurement that
exposed it.
\end{abstract}

\tableofcontents

\section{Introduction}

A genetic algorithm differs from a hill-climber in exactly one place: it
recombines. Everything else it does --- mutate, evaluate, select --- a
mutation-only search also does. The value of crossover is therefore the load
bearing claim of the whole family, and in evolutionary program synthesis it is
usually inherited from the schema-processing arguments of the string-genome era
rather than measured on programs.

This programme has approached the question three times from the detection side.
A quartet certificate refutes the four-point condition with an integer
comparison a proof kernel can recheck, and a recombinant genealogy violates that
condition where a clonal one does not; so if a program's behaviour were a
heritable sequence, the certificate would report whether crossover had
recombined. It is not. Behavioural profiles --- a program's outcome at each of a
battery of probes --- failed the test on real evolved populations in every
stratification we could construct, and a mutation-only arm, which is a tree by
construction, carried certificates at the same rate as a crossover arm. An
address-indexed ablation contribution, proposed as the heritable unit the
profiles lacked, transmitted at a median fraction of $0.200$--$0.545$ against a
pre-registered bar of $0.70$.

Each of those is a statement about \emph{measurability}. None of them says
whether crossover helps. That question requires running both arms to a verified
solution often enough for the difference to be visible, and the verification is
the expensive part.

\subsection{What this paper reports}

Two results, one methodological and one empirical, plus a mechanism and an
attempted repair.

\begin{enumerate}
\item \textbf{The verification cost was ours, not Lean's} (\S\ref{sec:atomicity}).
A compile containing a false obligation exits nonzero and still reports every
declaration's axiom dependencies. Our wrapper threw that away on the exit code.
Three rounds of isolation machinery were built on the resulting belief that a
batch is atomic; none of it was needed.

\item \textbf{Crossover significantly reduces the solve rate}
(\S\ref{sec:measurement}). $47.0\%$ against $29.0\%$ over $200$ paired seeds,
$p = 2.98\times10^{-4}$, with the evaluation count unchanged conditional on
success.

\item \textbf{The mechanism is the context coefficient}
(\S\ref{sec:mechanism}). An expression is affine in a subtree's value and the
coefficient depends on the operators above it, so a transplant preserves syntax
and not contribution.

\item \textbf{The repair this suggests fails, refuted by its own control}
(\S\ref{sec:homology}). Restricting recombination to sites with a matching
coefficient sign pattern does not recover mutation-only performance, and the
anti-homologous control is not worse than it.
\end{enumerate}

\subsection{What this paper does not claim}

The measurement is of one search: a closed grammar $\{+,-,*\}$ over four
variables and three rational constants, one calibrated target, one budget, one
steady-state population scheme with tournament selection, one crossover rate.
It is evidence about crossover in \emph{this} setting and not a general theorem
about recombination. The full specification target of the gate we use was solved
by neither arm at this budget, so nothing here is established at the size the
gate itself cares about. These limits are restated with the results rather than
deferred to a closing paragraph.

\section{The atomicity artifact}\label{sec:atomicity}

\subsection{What we believed, and recorded as measured}

Our verifier exposes a function that compiles a Lean file and returns one result
per declared theorem, where a declaration is clean only if the source passes a
forgery scan, the compile succeeds, and the kernel's \texttt{\#print axioms}
line for that declaration is free of \texttt{sorryAx} and of any axiom outside
the whitelist $\{\texttt{propext}, \texttt{Classical.choice},
\texttt{Quot.sound}\}$.

Batching several candidates into one file amortises a large fixed cost: a
candidate's own elaboration is about $0.14$\,s against about $5.23$\,s of fixed
per-invocation overhead, so roughly $97\%$ of a verification is paid whatever
the file contains. But we observed that a batch containing one false obligation
returned a single aggregate failure with no per-candidate verdict, and recorded
that as a property of Lean: \emph{a compile is atomic; one failing obligation
fails the whole file and yields no per-declaration rows}.

\subsection{What is actually the case}

It is a property of the wrapper. The function returns an aggregate failure the
instant the process exit code is nonzero, discarding output the kernel has
already produced. A single invocation on eight candidates with two false exits
$1$ and prints all eight lines:

\begin{quote}\small\ttfamily
'At0.At0\_m' depends on axioms: [propext, Classical.choice, Quot.sound]\\
'At2.At2\_m' depends on axioms: [propext, sorryAx, Classical.choice, Quot.sound]
\end{quote}

A failed tactic proof is admitted as \texttt{sorryAx} and reported by name,
which is exactly the condition the gate already rejects on.

\subsection{Why reading past the exit code is still sound}

Acceptance never rested on the exit code. It rests on three conditions that are
unchanged: the source-forgery scan, which rejects \texttt{\#exit} truncation and
\texttt{\#print} string spoofing; every declared theorem having its own axioms
line, so a declaration that never elaborated prints nothing and fails as
unchecked; and that line being free of \texttt{sorryAx} and of non-whitelisted
axioms. A nonzero exit now says only that some declaration failed, which the
per-name checks catch. Nothing is accepted that the stricter path would have
accepted; the report is widened, not the trust surface.

\subsection{The cost of the belief}

Three rounds of machinery, each internally sound, all downstream of the wrong
premise.

\begin{center}
\begin{tabular}{lrrrr}
\toprule
failures of 8 & bisection & flat fall-back & individual & kernel report \\
\midrule
0 & 1  & 1 & 8 & \textbf{1} \\
1 & 7  & 9 & 8 & \textbf{1} \\
2 & 11 & 9 & 8 & \textbf{1} \\
4 & 15 & 9 & 8 & \textbf{1} \\
8 & -- & 9 & 8 & \textbf{1} \\
\bottomrule
\end{tabular}
\end{center}

Bisection was built first, on the assumption that isolating $k$ failures in
$O(k\log n)$ invocations beats re-running $n$. At $97\%$ fixed cost every branch
re-pays a near-whole verification, so it is worse than the flat re-run it was
meant to improve on. The flat fall-back replaced it. Then a break-even law was
derived --- batching pays only while $(1-p)^b > 1/b$, giving a break-even
failure rate near $23\%$ at $b=8$ --- which made the high-failure regime an
evolutionary search occupies look like a barrier to the very experiment this
paper runs. With no isolation to charge for, the saving is $b$ at any failure
rate and the barrier does not exist.

In the shipped path, measured at $200$ seeds of the experiment below, all
distinct solutions were gated in one invocation with every culprit named and no
false candidate accepted at any failure rate.

\begin{remark}
The general form is worth stating because it is cheap to guard against. A
measured \emph{limit} reported by code one owns is a property of that code until
it is checked one layer below. The defect survived an existing guard --- that a
gate can be measuring one's own choice --- because the wrapper's output looked
like physics rather than like a choice.
\end{remark}

\section{Measuring crossover}\label{sec:measurement}

\subsection{Design}

A genome is an expression tree over the binary operators $\{+,-,*\}$ with leaves
drawn from the variables $h,w,q,p$ and the constants $1, 2, \tfrac12$. Fitness is
the total absolute exact error against a target at twelve rational sample points,
computed with arbitrary-precision rationals, so a fitness of zero means exact
agreement at those points. Agreement is not equality; that is what the kernel is
for.

The two layers are deliberately different. Fitness guides the search and costs no
proof-assistant time. Acceptance is the Lean kernel, applied to the rendered
candidate against the rendered target, proved by \texttt{ring} with its axiom
line checked. A wrong fitness function would waste search effort and could never
cause a wrong acceptance.

\subsection{Calibration, and a non-discriminating first attempt}

Our first attempt ran both arms against the full target and both solved $0/24$,
the best fitness reached anywhere being about $55.4$ with nothing below $10$.
That is a non-discriminating comparison: two arms failing equally say nothing
about recombination, and reporting it as evidence against crossover would have
been reading a result off an instrument that cannot separate the arms.

Difficulty was therefore calibrated on a ladder of reachable targets of
increasing size, with the rung chosen by a criterion fixed in advance: the
\emph{pooled} solve rate over both arms must lie in $[0.3, 0.7]$. Pooling is the
point --- selecting the rung at which crossover looked best would be choosing
the apparatus to suit the answer.

\begin{center}
\begin{tabular}{lrrrr}
\toprule
rung & size & pooled rate & mutation & crossover \\
\midrule
L1 & 5  & 0.750 & 13/16 & 11/16 \\
L2 & 7  & 0.969 & 16/16 & 15/16 \\
L3 & 9  & \textbf{0.438} & 9/16 & 5/16 \\
L4 & 11 & 0.125 & 1/16 & 3/16 \\
L5 & 13 & 0.000 & 0/16 & 0/16 \\
L6 & 17 & 0.000 & 0/16 & 0/16 \\
\bottomrule
\end{tabular}
\end{center}

Only L3 fell in the band. Note that L2 is \emph{easier} than L1 despite being
larger: $q + hp - q$ reduces to $hp$, so semantic difficulty and node count come
apart, which is itself a caution against reporting program-search difficulty by
size.

\subsection{Controls}

The target is asserted reachable and scoring exactly zero before any arm runs,
so a failure to solve is informative. An unreachable target --- one requiring
division by a variable, outside the grammar --- was solved $0$ times out of all
control runs; an arm that solved it would indicate a fault in fitness or in the
gate. Arms share seeds and initial populations.

\subsection{Result}

\begin{center}
\begin{tabular}{lrrr}
\toprule
 & solved & median evaluations & test \\
\midrule
mutation-only & $94/200$ \ ($47.0\%$) & $1286.5$ & \\
crossover     & $58/200$ \ ($29.0\%$) & $1269.0$ & Fisher $p = 2.98\times10^{-4}$ \\
\midrule
paired, both solved ($n=35$) & \multicolumn{2}{l}{ratio $0.959$, cheaper on $18$} & sign test $p = 1.000$ \\
\bottomrule
\end{tabular}
\end{center}

The two measures disagree and that is the result. Crossover significantly
reduces the probability of solving at all, while leaving the evaluation count of
successful runs essentially untouched. It does not make successful searches
faster; it makes success less likely.

\begin{remark}
At $48$ seeds the paired median ratio was $1.500$ over $10$ pairs, which would
have been reported as crossover being $50\%$ more expensive. At $200$ seeds and
$35$ pairs it is $0.959$ with a sign test of $p=1.000$. The small-sample paired
estimate did not replicate; the solve-rate difference, only $p = 0.147$ at $48$
seeds, is the measure that held. We record this because the first direction seen
in a paired statistic on ten pairs is not a result.
\end{remark}

\section{The mechanism}\label{sec:mechanism}

\begin{definition}[Context coefficient]
For an expression tree $t$ over $\{+,-,*\}$ and an address $a$ occurring once in
$t$, the value of $t$ is affine in the value $v$ of the subtree at $a$:
$t = A + Bv$, with $A$ and $B$ depending only on the rest of the tree and on the
evaluation point. We call $B$ the \emph{context coefficient} of $a$.
\end{definition}

$B$ is exactly computable by reverse accumulation: one forward pass for every
node's value, then one backward pass seeded with $B(\text{root}) = 1$, where $+$
sends $B$ to both children, $-$ sends $B$ left and $-B$ right, and $*$ sends
$B\cdot\text{value}(\text{sibling})$ to each child. This is $O(|t|)$ per sample
point in exact rational arithmetic. We verified it against the defining two
evaluations $B = t|_{v=1} - t|_{v=0}$ on random trees, with exact agreement over
more than five hundred address--point pairs.

The diagnosis follows immediately. A donor subtree contributing
$B_{\text{donor}}\,v$ inside its own parent contributes $B_{\text{site}}\,v$
after transplant. Unless the two coefficients agree, the donor's
\emph{contribution} is rescaled or sign-inverted even though its syntax arrived
intact. Subtree crossover transmits structure and not function, which is what it
means to call it a macromutation. This agrees with an independent finding on
different evidence: transfer fidelity of a donor's ablation signature falls with
the number of multiplicative operators on the root-to-insertion path.

\begin{remark}
The naive implementation of $B$ --- two whole-tree evaluations per address ---
costs $O(|t|^2)$ per point and measured at $5.5\times$ the tree-point
evaluations of a plain arm. Had the homology arms below been compared on fitness
evaluations alone, that hidden cost would have made any win meaningless. Both
cost metrics are reported for every arm.
\end{remark}

\section{Homologous recombination, and why the repair fails}\label{sec:homology}

Biology does not recombine at arbitrary positions; it recombines at homologous
ones, and non-homologous rejoining is a damage pathway rather than a creative
one. Random subtree crossover is the non-homologous analogue. The diagnosis of
\S\ref{sec:mechanism} suggests the corresponding repair: admit a donor only into
a site that preserves its contribution, operationalised as requiring the sign
pattern of $B$ across the sample points to match the donor's.

\subsection{The prediction, registered before the run}

Recorded in the artifact before any arm executed: homology-restricted
recombination should raise the solve rate above random crossover's $29.0\%$,
ideally to at least mutation-only's $47.0\%$; anti-homologous recombination,
restricted to \emph{inverted} sign patterns, should do no better than random. We
stated the falsifier explicitly --- \emph{if the two homology arms behave alike,
the signature carries nothing and the mechanism is not what was claimed} --- and
counted fallbacks to mutation, taken when no admissible site exists, so that an
arm quietly becoming mutation-only could not read as a success of recombination.

\subsection{Result}

\begin{center}
\begin{tabular}{lrrr}
\toprule
arm & solved & rate & median tree-point evaluations \\
\midrule
mutation-only              & $94/200$ & $0.470$ & $48\,000$ \\
crossover (random)         & $58/200$ & $0.290$ & $48\,000$ \\
homologous crossover       & $70/200$ & $0.350$ & $94\,872$ \\
anti-homologous (control)  & $79/200$ & $0.395$ & $94\,812$ \\
\bottomrule
\end{tabular}
\end{center}

\begin{center}
\begin{tabular}{llr}
\toprule
comparison & what it tests & Fisher exact $p$ \\
\midrule
homologous vs.\ random crossover & does homology beat random? & $0.238$ \\
\textbf{anti vs.\ homologous} & \textbf{the control must be worse} & $\mathbf{0.408}$ \\
mutation-only vs.\ homologous & does the repair recover mutation? & $\mathbf{0.019}$ \\
mutation-only vs.\ anti & & $0.158$ \\
\bottomrule
\end{tabular}
\end{center}

The repair fails on every clause of its own prediction. Homology restriction
does not significantly beat random crossover. It does not recover mutation-only,
which beats it significantly. And the negative control --- built to be worse ---
is nominally \emph{higher} and statistically indistinguishable from the
treatment, which is precisely the outcome we had registered in advance as
meaning the coefficient sign pattern does not carry the mechanism.

\subsection{Ruling out the obvious confound}

Both homology arms fall back to mutation when no admissible site exists, so a
sceptic's first move is to attribute their edge over random crossover to extra
mutation. The counts exclude it. Fallback rates are $2.35\%$ (homologous) and
$3.91\%$ (anti-homologous) of recombination attempts; since recombination is
attempted on half of iterations, the effective mutation fractions are $0.5117$
and $0.5196$ against random crossover's $0.5000$. A difference under two
percentage points in the operator mix cannot account for a $10.5$ percentage
point spread in solve rate. The ordering of the three crossover arms is
therefore \emph{not} explained by the fallback --- and we leave it unexplained
rather than supply a story, since no adjacent pair among them is significant.

\subsection{What survives}

What survives is the negative result, now from five independent directions:
behavioural profiles are not heritable sequences; the address-indexed ablation
contribution is not a site-wise heritable character; a mutation-only arm carries
quartet certificates at the same rate as a crossover arm; random subtree
crossover significantly lowers the solve rate; and the homology repair that the
exact mechanism recommends does not recover the loss, with its own control
refuting the mechanism. In this representation, the only operator that reliably
works is mutation.

We note what this does \emph{not} license. It does not show that no
contribution-preserving recombination can work --- only that this one, defined by
the sign pattern of the context coefficient, does not, and that its failure
cannot be rescued by the mechanism we proposed. A coefficient-matching scheme
that \emph{rescales} the donor to preserve $B\cdot v$ rather than merely matching
the sign of $B$ remains untested, and is the obvious next candidate. We report
the failure now rather than after finding a variant that works, because the
prediction and its falsifier were fixed before the run.

\section{Preserving the contribution exactly}\label{sec:preserving}

Edition~1 closed by naming the successor it had not tested: matching
$\operatorname{sign}(B)$ failed, so the obvious candidate was to \emph{rescale}
the donor and preserve $B\cdot v$ itself. Edition~2 tests it.

\subsection{The operator}

A site is admissible only when $B_{\text{donor}} = k\,B_{\text{site}}$ for a
\emph{single} rational $k$ at every sample point --- one scalar cannot correct a
point-varying ratio --- and the donor is then inserted as $k\cdot D$, restoring
$B_{\text{donor}}\,v$ exactly. The case $k=1$ needs no wrapper and is common,
since $B=1$ at every address reached through only $+$ and the left of $-$, which
is also why an earlier fidelity measurement found $1.000$ in the $\{+,-\}$ arm.

The invariant is measured, not asserted: $1614$ of $1614$ contribution
comparisons are exactly equal before and after the graft, in exact rational
arithmetic, with admissible sites found in $269$ of $400$ random donor/site
pairs, so the condition is neither vacuous nor trivial.

\subsection{A sharper control than edition 1's}

\texttt{same\_sites\_unscaled} selects the \emph{same} constant-ratio sites and
omits the scaling. This separates the two things the operator does --- choosing a
site, and correcting the contribution once there --- which edition~1's design
could not.

\subsection{Result}

\begin{center}
\begin{tabular}{lrrr}
\toprule
arm & solved & rate & median evaluations \\
\midrule
mutation-only                     & $94/200$ & $0.470$ & $1286.5$ \\
crossover (random)                & $58/200$ & $0.290$ & $1269.0$ \\
homologous (edition 1)            & $70/200$ & $0.350$ & $1386.5$ \\
anti-homologous (edition 1)       & $79/200$ & $0.395$ & $1712$ \\
\textbf{contribution-preserving}  & $73/200$ & $0.365$ & $1011$ \\
same sites, unscaled (control)    & $52/200$ & $0.260$ & $1358.5$ \\
\bottomrule
\end{tabular}
\end{center}

\begin{center}
\begin{tabular}{llrl}
\toprule
comparison & role & $p$ & Bonferroni \\
\midrule
preserving vs.\ unscaled & isolating test & $0.031$ & no \\
preserving vs.\ mutation-only & does it rescue crossover? & $0.042$ & no \\
\textbf{preserving vs.\ crossover} & \textbf{registered falsifier} & $\mathbf{0.136}$ & no \\
unscaled vs.\ crossover & sanity & $0.576$ & no \\
preserving vs.\ homologous & exact vs.\ sign & $0.835$ & no \\
\bottomrule
\end{tabular}
\end{center}

\subsection{The falsifier fired}

We registered, before running: \emph{if contribution-preserving does not
significantly exceed random crossover, then exact contribution preservation is
also not the missing ingredient, and the contribution-transfer story is refuted
as an explanation rather than merely in its sign-matching form.} It does not
($p = 0.136$). The hypothesis fails the bar we set for it.

Five comparisons were run, so a family-wise rate of $0.05$ requires
$\alpha = 0.01$, and \emph{nothing survives correction} --- including the two
raw-significant results.

\subsection{What the data do support, at their real strength}

The isolating comparison is the interesting one. Holding site selection fixed,
turning the scaling on beats turning it off at $p = 0.031$, and the sanity check
holds: the unscaled arm is indistinguishable from random crossover
($p = 0.576$), so site selection alone buys nothing and any effect lives in the
scaling. That is the first positive mechanistic signal in this programme.

It is also not enough, and we decline to promote it. It fails correction for the
five comparisons, and it is the comparison against the \emph{weakest} arm;
elevating it to the headline would be choosing the test after seeing the data,
which is the defect this programme keeps finding in its own work. Three further
things are therefore \emph{not} claimed: that the scaling works, since the
registered bar was exceeding random crossover and it was not met; that the
operator is faster conditional on solving, though its median of $1011$
evaluations against mutation-only's $1286.5$ is suggestive at $p = 0.057$; and
that exact preservation beats sign-matching, since $p = 0.835$ --- the stronger
invariant, verified exact at every one of $1614$ comparisons, bought nothing
measurable over matching a sign.

A $6$-seed smoke test had shown $3/6$ at a median of $475$ evaluations. It was
reported as a reason to run at $200$ seeds rather than as a result, and at $200$
the arm fails its bar. That is the third small-sample estimate in this work not
to replicate.

\subsection{What this leaves}

Six independent results agree that subtree recombination does not help in this
representation, and \emph{both} mechanism-targeted repairs --- matching
$\operatorname{sign}(B)$, and preserving $B\cdot v$ exactly --- fail their
pre-registered bars. The contribution-transfer account of why crossover fails
survives only as a hint from one uncorrected sub-comparison, which is not enough
to call it the cause. We state that rather than the alternative, which would be
to keep proposing variants until one clears a bar set afterwards.

\section{Reproduction}

Every numeric claim above is produced by a script in the accompanying
repository. The search, calibration, controls and gate are a single file; the
verifier changes and their tests are in the main project. The artifacts carry
the measured values, the pre-registered criteria, and the limits of each claim,
including the ones that qualify it.

\section{Related work and honest placement}

This is a negative result about an operator, measured in one setting, with a
mechanism and a failed or successful repair reported as it came out. It does not
overturn the genetic-algorithm literature, where crossover's value has been both
defended and attacked for decades and where representation is known to matter
enormously. Its contribution is methodological as much as empirical: the
measurement was affordable only after we found that our own verifier had been
reporting a limit of its own making, and the discipline that caught that ---
measuring a limit one layer below the code that reports it --- is the part most
likely to transfer.

\end{document}
